\scanpage{6}
\diagram{\begin{tikzcd}[row sep=small,column sep=large]
 & \pi \arrow[d,"u"'] & & & \\
 & G(\mathbf Q_p) \arrow[d,phantom,"\unsure{\bigcap}" description] & & & \\
 & G(\mathbf C) \arrow[d,phantom,"\unsure{\alpha|\simeq}" description] & & & \\
1 \arrow[r,no head] & G'(\mathbf C) \arrow[r] \arrow[d,phantom,"\unsure{\bigcup}" description] & E \arrow[r] & \pi \arrow[l,bend right=35,"\sigma_1"'] \arrow[l,bend left=35,"\sigma_2"] \arrow[r,no head] & 1 \\
 & G'(\mathbf Q_p) & & & \\
 & G_m(\mathbf Q_p)=\mathbf Q_p^* \arrow[u,hook,"j_2"] & & & \\
 & \pi \arrow[u,"\chi"] & & &
\end{tikzcd}}

\deleted{\unsure{on}} $\alpha$ permet d'identifier une extension cat. \unsure{en} $E$ de $\pi$ par $G'(\mathbf C)$ \unsure{en} une \unsure{ext\textsuperscript{on}} cat. de $\pi$ par $G(\mathbf C)$.

Le splitting $\sigma_2$ de la première est tel que $\sigma_2(\pi)$ commute à $G'(\mathbf Q_p)$ et a fortiori à $j_2\chi(\pi)$. On trouve donc
\[
\left\{\begin{array}{l}
\gamma\longmapsto\widetilde\sigma_2(\gamma)=\sigma_2(\gamma).j_2\chi(\gamma)\\
\pi\longrightarrow E
\end{array}\right.
\]
un autre splitting.

Le splitting $\sigma_1$ de la seconde commute à $G(\mathbf Q_p)$, et à fortiori à $u(\pi)$. On trouve donc
\[
\gamma\longmapsto\sigma_1(\gamma)u(\gamma)=\widetilde\sigma_1(\gamma).
\]

On écrit que $\widetilde\sigma_1(\gamma)=\widetilde\sigma_2(\gamma)$

i.e.
\[
\sigma_2(\gamma)j_2(\chi(\gamma))=\sigma_1(\gamma)u(\gamma)
\]
donc
\[
\underbrace{\sigma_1(\gamma)^{-1}\sigma_2(\gamma)}_{\varphi(\gamma)}
=\deleted{\unsure{j_2}}\underset{\substack{\rotatebox{90}{$\in$}\\G(\mathbf Q_p)}}{u(\gamma)}.
\underset{\substack{\rotatebox{90}{$\in$}\\j_2(\mathbf Q_p^*)}}{\underbrace{j_2(\chi(\gamma^{-1}))}}
\]

1°) Sur $\operatorname{Ker}\chi$, on a $\varphi(\gamma)=u(\gamma)$. donc $u$ connu \unsure{sur} \deleted{$\unsure{\operatorname{Ker}\chi}$}

2°) $\operatorname{Ker}u=\deleted{\chi(\gamma)}\{\gamma,\chi(\gamma)=1,\varphi(\gamma)=1\}$ connu. Soit $\unsure{\Pi_1|\Pi'_1}=\pi/\operatorname{Ker}u$. C'est un groupe analytique $p$-adique, et a
\[
\Pi_1\xrightarrow{\chi_1}\mathbf Q_p^*\quad\text{induit par }\unsure{\rho}.
\]
\margin{5}
